\documentclass[11pt,a4paper]{article}

\usepackage[margin=1in]{geometry}
\usepackage{amsmath,amssymb}
\usepackage{physics}
\usepackage{booktabs}
\usepackage{hyperref}
\usepackage{enumitem}

\newenvironment{resultbox}[1][Result]{%
  \begin{quote}\noindent\textbf{#1.}\ }{\end{quote}}
\newenvironment{gapbox}[1][Open problem]{%
  \begin{quote}\noindent\textbf{#1.}\ }{\end{quote}}

\newcommand{\half}{\tfrac{1}{2}}
\newcommand{\mzero}{m_0}
\newcommand{\seriesname}{Saturated Superfluid Vacuum}
\newcommand{\papernumber}{Goldstone}

\title{\textbf{\seriesname~\papernumber}\\[0.4em]
The Photon as a Goldstone Mode and the Chiral-Shear Near Field}
\author{Stig Norland \\ Independent Researcher \\ Bergen, Norway}
\date{\today}

\begin{document}
\maketitle

\begin{abstract}
This paper isolates the carrier problem in the electromagnetic sector of the
Saturated Superfluid Vacuum (SSV) series.  Static Coulomb and magnetic
near-field structure arises from chiral-shear flow gradients around toroidal
charged vortices, with coupling strength set by the dimensionless parameter
$\alpha$.  Identifying the propagating photon with a chiral-shear mode does not
work: linearising the full LogSE + chiral-shear + CP$^1$ functional around the
uniform vacuum shows the chiral-shear term is \emph{silent} in the linear
spectrum --- the linear-order current is a pure gradient and the
$|\nabla\times\mathbf{j}|^2$ energy is quartic --- so no propagating
chiral-shear mode exists at any speed, and the internal-direction channel
disperses as $\omega=\hbar k^2/2m_0$: gapless, quadratic, charge-decoupled at
linear order.

Eliminating that candidate does not establish the alternative.  The saturated
order parameter has a longitudinal Goldstone/Bogoliubov sound mode at $c$, but
a scalar phase mode alone is not a photon: it supplies neither two transverse
polarizations nor Maxwell gauge structure, and the CP$^1$ direction is
charge-decoupled at linear order, so its transport cannot simply be asserted to
provide them.  The chiral-shear sector remains physically real and is assigned
to static and near-field structure --- Coulomb pressure, magnetic circulation,
Aharonov--Bohm phase, vortex-ring self-energy.

The task is therefore a two-channel derivation: obtain the radiative Maxwell
sector from the Goldstone phase channel at speed $c$ while preserving the
static Coulomb and magnetic results as the near-field limit of chiral-shear
topology.  That derivation is \textbf{not} claimed here.  This paper states the
consistency conditions, identifies the matching problem between the channels,
and fixes the success criterion: a transverse, gauge-like two-polarization
sector with displacement-current coefficient $1/c^2$.
\end{abstract}

\newpage
\tableofcontents
\newpage

% ---------------------------------------------------------------------
\section{Introduction}
\label{sec:introduction}
% ---------------------------------------------------------------------

The electromagnetic sector of SSV has split into two logically distinct
pieces:
\begin{enumerate}[label=(\roman*),leftmargin=2em]
    \item a \emph{static near-field sector}, where charge, Coulomb force,
          magnetic circulation, and AB phase arise from the
          topology and chiral-shear flow of toroidal vortices;
    \item a \emph{radiative sector}, where the observed photon must
          propagate at the measured vacuum speed $c$.
\end{enumerate}
Paper~II developed the first piece but used the wrong carrier for the
second.  It treated the chiral-shear mode as the photon and assigned it the
speed $c_\perp$.  This cannot be the final physical interpretation because
the observed photon propagates at $c$.

Analogue-gravity and superfluid-vacuum literature motivates emergent
collective fields~\cite{volovik2003}, while
Zloshchastiev~\cite{zloshchastiev2020} supplies the logarithmic wave
equation.  Neither cited source establishes the SSV-specific identification
of the physical photon with the scalar phase mode of $\Psi$.  That phase mode
propagates at
the medium sound speed,
\begin{equation}
    c_s(\rho_0)=c,
    \label{eq:cs_c}
\end{equation}
and is therefore compatible with observed light propagation and with
gravitational-wave constraints such as GW170817.

The question is not simply whether one can rename the photon.  The hard
question is whether the successful static electromagnetic mechanisms of
Paper~II survive the replacement:
\[
    \hbox{chiral-shear photon at }c_\perp
    \quad\longrightarrow\quad
    \hbox{Goldstone photon at }c .
\]
The present paper formulates that problem.

\begin{resultbox}[Candidate programme]
The Goldstone/Bogoliubov mode supplies a candidate carrier speed $c$, while
the chiral-shear field supplies a candidate static near-field sector.  A
photon identification requires a derivation of two transverse polarizations,
Maxwell dynamics and source matching; none follows from the scalar sound mode
alone.
\end{resultbox}

% ---------------------------------------------------------------------
\section{The Problem in Paper II}
\label{sec:problem}
% ---------------------------------------------------------------------

Paper~II recovered a set of electromagnetic structures from chiral vortex
topology:
\begin{align}
    F_C(r) &= \frac{\alpha\hbar c}{r^2}, \label{eq:coulomb}\\
    \mathbf{B} &\sim \nabla\times\mathbf{v}_\perp, \label{eq:Bcurl}\\
    \gamma_{\rm AB} &= \frac{q}{\hbar}\oint\mathbf{A}\cdot d\mathbf{l}.
    \label{eq:AB}
\end{align}
Here $\mathbf{v}_\perp$ is the chiral-shear flow associated with the
toroidal vortex, and $\alpha$ is the empirical dimensionless strength of
that sector.

The same section then tentatively identified the propagating photon with a
transverse chiral-shear phonon satisfying
\begin{equation}
    \partial_t^2\delta\theta
    =
    c_\perp^2\nabla^2\delta\theta,
    \qquad
    c_\perp \neq c.
    \label{eq:old_photon}
\end{equation}
This produces the speed problem.  If the physical photon propagated at
$c_\perp$, vacuum light speed would be $c_\perp$, not $c$.  The earlier
patch that $\alpha\to1$ in free vacuum is not available inside SSV because
vacuum \emph{is} the saturated medium.

The diagnosis is therefore:
\begin{equation}
    \hbox{static chiral-shear mechanism: keep,}
    \qquad
    \hbox{radiative photon identification: replace.}
    \label{eq:diagnosis}
\end{equation}

The SSV issue~138 computation subsequently showed that
equation~\eqref{eq:old_photon} was not merely the wrong photon --- it does
not arise from the functional at all: around the uniform vacuum the
chiral-shear term contributes nothing at linear order, and the
internal-direction channel it was thought to describe is first-order in
time with $\omega=\hbar k^2/2m_0$.  The ``slow second light'' has no
linear mode behind it; the diagnosis above is thereby upgraded from an
audit conclusion to a computed result.

% ---------------------------------------------------------------------
\section{Goldstone Channel of the Saturated Medium}
\label{sec:goldstone}
% ---------------------------------------------------------------------

Write the order parameter as
\begin{equation}
    \Psi(\mathbf{x},t)
    =
    \left(\sqrt{\rho_0}
    +\frac{\delta\rho}{2\sqrt{\rho_0}}\right)
    e^{i\theta_0+i\delta\theta}.
    \label{eq:psi_expand}
\end{equation}
Linearisation gives coupled density and phase perturbations:
\begin{itemize}[leftmargin=2em]
    \item $\delta\rho$ is the density component;
    \item $\delta\theta$ is the phase component.
\end{itemize}
Together they form the longitudinal Bogoliubov sound branch.  At long
wavelengths its phase component obeys
\begin{equation}
    \partial_t^2\delta\theta
    =
    c^2\nabla^2\delta\theta,
    \label{eq:goldstone_wave}
\end{equation}
with dispersion
\begin{equation}
    \omega = ck .
    \label{eq:goldstone_disp}
\end{equation}

This branch has the required propagation speed, but it is a scalar
longitudinal sound mode.  Speed agreement is necessary and not sufficient
for identifying it with physical electromagnetic radiation.

\begin{resultbox}[Candidate channel assignment]
The Goldstone/Bogoliubov branch is the only current candidate carrier at
$c$.  The chiral-shear channel is a candidate near-field structure of
charged defects.  Their electromagnetic interpretation remains unproved.
\end{resultbox}

% ---------------------------------------------------------------------
\section{Two-Channel Electromagnetism}
\label{sec:two_channel}
% ---------------------------------------------------------------------

The corrected EM sector should be written as a two-channel theory:
\begin{align}
    \hbox{near field:}\qquad
    &\mathbf{v}_\perp,\ \Phi_\perp,\ \mathbf{A}_\perp,
    &&\hbox{chiral-shear/topological sector},\\
    \hbox{radiation:}\qquad
    &\delta\theta_{\rm G},
    &&\hbox{Goldstone phase channel}.
    \label{eq:two_channel}
\end{align}
The near field knows about charge chirality and the dimensionless coupling
$\alpha$.  The radiative field knows about propagation at $c$.

This is not unusual from the point of view of ordinary electrodynamics.  In
QED, the Coulomb field is not a freely propagating transverse photon; it is
the constrained or longitudinal part of the gauge field.  Radiation is
carried by transverse on-shell photon modes.  SSV should reproduce this
division mechanically:
\begin{align}
    \hbox{Coulomb interaction}
    &\leftrightarrow
    \hbox{static chiral-shear pressure constraint},\\
    \hbox{on-shell light}
    &\leftrightarrow
    \hbox{Goldstone phase wave at }c.
    \label{eq:qed_analogy}
\end{align}

The task is to derive the effective Maxwell system in which the source
sector is chiral-shear/topological, while the homogeneous radiative sector
propagates at $c$:
\begin{align}
    \nabla\cdot\mathbf{E} &= \rho_{\rm ch}/\epsilon_0,\\
    \nabla\cdot\mathbf{B} &= 0,\\
    \nabla\times\mathbf{E} &= -\partial_t\mathbf{B},\\
    \nabla\times\mathbf{B} &= \mu_0\mathbf{J}
       +\frac{1}{c^2}\partial_t\mathbf{E}.
    \label{eq:maxwell_target}
\end{align}
Relative to the provisional Paper~II derivation, the displacement-current
coefficient must be $1/c^2$ rather than $1/c_\perp^2$.

% ---------------------------------------------------------------------
\section{What Must Survive}
\label{sec:survive}
% ---------------------------------------------------------------------

The Goldstone correction is acceptable only if it preserves the successful
static and topological content of Paper~II.  The following structures must
survive.

\paragraph{Coulomb strength.}
The static force must remain
\begin{equation}
    F_C(r)=\frac{\alpha\hbar c}{r^2}.
    \label{eq:coulomb_survive}
\end{equation}
In the corrected reading, $\alpha$ is not a photon-speed ratio.  It is the
strength of the chiral-shear near-field coupling.

\paragraph{Magnetic topology.}
The magnetic field must remain the curl of a vector flow,
\begin{equation}
    \mathbf{B}\propto\nabla\times\mathbf{v}_\perp,
    \label{eq:B_survive}
\end{equation}
so that $\nabla\cdot\mathbf{B}=0$ remains a vector identity and magnetic
monopoles remain topologically forbidden.

\paragraph{Aharonov--Bohm phase.}
The phase shift must remain a Berry phase of the charged ring in a
background circulation:
\begin{equation}
    \gamma_{\rm AB}
    =
    \frac{q}{\hbar}\oint\mathbf{A}\cdot d\mathbf{l}.
    \label{eq:AB_survive}
\end{equation}
This is naturally near-field/topological rather than radiative.

\paragraph{Anomalous magnetic moment.}
The one-loop self-energy picture must be reinterpreted carefully.  The
virtual exchange responsible for $g-2$ may continue to involve the
chiral-shear near field, but the on-shell photon propagator must belong to
the Goldstone channel.  The calculation therefore has to distinguish
virtual constrained exchange from physical radiative quanta.

\begin{gapbox}[Open matching problem]
The central technical challenge is to show that the chiral-shear near field
and the Goldstone radiation field are not two unrelated electromagnetic
fields, but two limits of one effective EM response of the saturated medium.
\end{gapbox}

% ---------------------------------------------------------------------
\section{The Matching Calculation}
\label{sec:matching}
% ---------------------------------------------------------------------

A complete derivation should start from the full SSV functional,
\begin{equation}
    S[\Psi]
    =
    \int dt\,d^3x\,
    \left[
        \mathcal{L}_{\rm LogSE}(\Psi)
        +
        \mathcal{L}_{\perp}(\Psi,\nabla\Psi)
        +
        \mathcal{L}_{\rm defect}
    \right],
    \label{eq:full_action}
\end{equation}
expand around a charged toroidal background $\Psi_{\rm e}$, and separate
fluctuations into constrained near-field and propagating Goldstone parts:
\begin{equation}
    \Psi
    =
    \Psi_{\rm e}
    + \delta\Psi_{\rm near}
    + \delta\Psi_{\rm rad}.
    \label{eq:split}
\end{equation}
The desired effective action has the schematic form
\begin{equation}
    S_{\rm eff}
    =
    \int dt\,d^3x\,
    \left[
      \frac{\epsilon_0}{2}E^2
      -\frac{1}{2\mu_0}B^2
      + J^\mu A_\mu
    \right],
    \label{eq:Seff}
\end{equation}
with two requirements:
\begin{align}
    \hbox{static limit:}\quad
    &\nabla^2\Phi = -\rho_{\rm ch}/\epsilon_0,
    &&F_C=\alpha\hbar c/r^2,\\
    \hbox{radiative limit:}\quad
    &\partial_t^2 A_i - c^2\nabla^2 A_i = 0,
    &&\omega=ck.
    \label{eq:limits}
\end{align}

The calculation is therefore a matched asymptotic problem.  It must test
whether chiral-shear topology can supply the source and static response and
whether a far-field disturbance can project onto a transverse sector
transported at the Goldstone-branch speed.

% ---------------------------------------------------------------------
\section{Avoiding New Circularity}
\label{sec:circularity}
% ---------------------------------------------------------------------

The Goldstone photon correction must not simply insert Maxwell theory by
hand.  The following shortcuts would not count as derivations:
\begin{enumerate}[label=(\roman*)]
    \item postulating a separate electromagnetic gauge field in addition to
          $\Psi$;
    \item replacing $c_\perp$ by $c$ in Paper~II's equations without
          deriving the channel projection;
    \item keeping Coulomb from chiral shear and radiation from Goldstone
          without deriving their matching at sources;
    \item using observed light speed to tune a coefficient in the
          chiral-shear sector.
\end{enumerate}

A valid derivation should show that the same saturated order parameter
supports both the constrained static response and the propagating radiation
response, with no additional gauge field inserted by hand.

\begin{resultbox}[Success criterion]
The Goldstone photon programme succeeds if the full SSV linearisation around
charged toroidal defects yields Maxwell electrodynamics with static Coulomb
strength $\alpha\hbar c$, magnetic topology from chiral-shear circulation,
and on-shell radiation satisfying $\omega=ck$.
\end{resultbox}

% ---------------------------------------------------------------------
\section{Consequences}
\label{sec:consequences}
% ---------------------------------------------------------------------

If the matching calculation succeeds, several tensions in Paper~II are
resolved at once.

\paragraph{Light and gravity share the correct speed.}
Both electromagnetic radiation and gravitational disturbances can be read
as Goldstone/Bogoliubov-channel excitations at $c$, consistent with
GW170817.

\paragraph{\texorpdfstring{$\alpha$}{alpha} is not a speed ratio.}
The fine-structure constant becomes a chiral-shear/topological coupling and
defect-geometry parameter, not a ratio involving the observed photon speed.
This matches the separate alpha-paper proposal
$\alpha=\xi/R^*$.

\paragraph{The old EM derivation becomes a near-field derivation.}
The Paper~II chiral-shear equations are not discarded; they are reclassified.
They describe constrained near-field structure and virtual exchange rather
than the on-shell photon.

\paragraph{Magnetic monopoles remain forbidden.}
Because the magnetic field still descends from a curl of the chiral-shear
flow, $\nabla\cdot\mathbf{B}=0$ remains topological rather than postulated.

% ---------------------------------------------------------------------
\section{Discussion}
\label{sec:discussion}
% ---------------------------------------------------------------------

The Goldstone photon problem is a cleanup calculation, but it is not a
minor one.  It decides whether the SSV electromagnetic sector can be made
internally coherent after the propagation-speed audit.  The desired result
is conceptually simple: sources live in chiral topology; radiation lives in
the Goldstone phase channel.  The hard part is proving that these are two
faces of one effective electromagnetic field.

This also clarifies the status of Paper~II.  Its static Coulomb and
magnetic results remain valuable because they identify how charge and
magnetic topology arise mechanically.  Its original radiative photon
identification should be treated as provisional scaffolding.  The present
paper states the corrected target.

% ---------------------------------------------------------------------
\section{Conclusion}
\label{sec:conclusion}
% ---------------------------------------------------------------------

The current SSV functional rules out the former slow chiral-shear photon.
Its remaining Goldstone/Bogoliubov sound branch propagates at $c$, but the
paper has not derived a photon from it.  The chiral-shear near-field
interpretation is likewise conditional on the matching calculation.

The open derivation is the matching calculation:
\[
    \boxed{
    \hbox{chiral-shear near field}
    \quad\longleftrightarrow\quad
    \hbox{Goldstone radiation at }c
    }
\]
If this can be derived from the full SSV functional without adding a
separate gauge field, the electromagnetic sector becomes substantially more
coherent.  If it cannot, the SSV force-sector programme must be revised at
one of its central joints.

\bibliographystyle{unsrt}
\bibliography{../cited/references}

\end{document}
